\subsection{Workflow guidance, scientific judgment, and outlook}

Tool and task graphs organize computational execution. Cloudra's Delivery DAG constrains the mechanism representation submitted by the agents: nodes represent species occurrences and edges represent elementary steps, carrying molecular geometries, bonding patterns, and reaction connectivity. Cheminformatics algorithms establish and compare molecular identities and atom correspondences; frequency and IRC/endpoint calculations determine which stationary points and connections the calculations support. The 50-task comparison measures end-to-end mechanism-calculation performance with these methods. Guided had higher mean self-consistency and accuracy than Standard; in Hard tasks, their weighted contributions to the gain were +0.1353 and +0.1274. HH1 and HI1 report reference-consistent barriers; identifying the minima connected by those saddles still requires IRC and endpoint calculations.

HA8 and NU6 show two assignments that must be checked after geometry optimization. In HA8, the stable reactant minimum defines the barrier reference, while the IRC connects the saddle to a local reactive conformer. In NU6, an initial guess for betaine-to-product closure repeatedly optimized toward the direct reactant-to-product saddle region. The resulting saddle must therefore be assigned from its optimized structure and path. The DAG records the reactant conformer, saddle, IRC endpoints, and assigned step so that the agent can inspect which structures each energy belongs to after a new calculation. This use of recorded calculation results is consistent with the role of external planning and verification feedback in language-model workflows.\citep{valmeekam2023planning}

A missing endpoint optimization calls for another endpoint calculation; a saddle that repeatedly returns to another route calls for a revised guess or pathway hypothesis. An electronic state beyond the method's scope requires reassessing the method and the requested conclusion. The agent must compare its reaction description with the computed structures and paths, because a plausible explanation can persist after the underlying result has changed.\citep{turpin2023unfaithful} Record retained the tools and obtained lower self-consistency than Guided, particularly on Normal and Hard tasks. The Brief-errors comparison also differed by task: detailed feedback accompanied higher Guided means on Normal and Hard, while Brief-errors had the higher Impossible mean. These results support examining whether feedback directs the agent to the relevant geometry, calculation, or method limitation, consistent with studies of tool feedback in reasoning correction.\citep{ref18,huang2024selfcorrect}

Brief-errors scored higher on IA1 and IA9, whereas Guided scored higher on IG8. These crossed outcomes motivate a task-dependent hypothesis: detailed diagnostics may help repair a TS or IRC calculation, with limited benefit from further local repair when the method or target needs reconsideration. The present comparisons document final task outcomes; the causal effect of diagnostic length and its underlying internal mechanism remain questions for dedicated study.

In G2, checks on a companion edge blocked numerical declarations, after which the agent removed values already supported by native calculations. Supervisor feedback restored the fields; further rejection preceded their deletion from the final graph. While the pathway is being tested, the numerical observation and its native log should remain available, with the current validation status specifying which mechanistic conclusions the value supports. G2 retained partial self-consistency credit and lost accuracy credit when the numerical fields were removed.

P2 contains one retained run per task and configuration. Its task-cluster bootstrap intervals describe variation across the sampled task mix; same-task stochastic variation requires independent repeats. The separate P1b study supplies three repeats on each of four tasks under each configuration; these 48 runs cover four of the 50 P2 tasks. The complete organized P2 table covers all 200 runs and their three score components. Some Hard-task accuracy targets concern qualitative channel identity, so the reported gain includes pathway identification and evidence completion as well as numerical barriers. Repeating the full task set would characterize run-to-run variability and the frequency with which each research route is selected.

The Impossible tier concerns targets judged physically nonexistent under the task's permitted methods; these judgments are specific to that method description. Its component values also use negative-target rules: self-consistency 1 on an explicit negative conclusion is distinct from a positive TS--IRC--endpoint chain. In the 10 Impossible cases with self-consistency 1 and accuracy 0, the score must be read alongside the particular state, pathway, or property that the run characterized or rejected.

The 12 pilot batches measure real elapsed time for four tasks at concurrency two, including the original failures and stopping outcomes. The P2 effort archive covers selected deep-dive studies, with 15 tasks shared across configurations; cumulative job-hours represent recorded quantum-chemical process durations, with a documented native-log reconstruction for Standard/M1. They are distinct from core-hours, financial cost, and end-to-end P2 research latency. Concurrent calculations shared computational resources, and CPU queueing may have delayed individual results. Separating queueing, calculation, and agent-active time would show how much of each run was spent waiting, computing, or deciding what to do next.

Standard's missing graphs followed premature completion while calculations were still pending; the Guided/HI1 session remained waiting after the completion notification. In both cases, a completed calculation needed to resume the unfinished study before the final report was assembled. A further search problem is repeated optimization toward the same stationary-point region, as in NU6. Candidate-selection methods could compare new saddles with earlier attempts and help the agent choose the next guess or pathway hypothesis.

Cloudra's default preorganization samples 32 rigid relative placements, applies local MMFF94 relaxation, and selects candidates by energy and clustering. Frequency verification characterizes the candidate stationary point; a search for lower conformers requires additional sampling. HA8 requires both the specified stable reference and the local reactive conformer; a global-minimum claim additionally requires an assessment of conformer coverage. CREST is available through a separate optional backend; its coverage across the 200 runs remains to be established. Highly flexible systems may require broader sampling and more quantum-chemical calculations. Further work should measure which low-energy conformers are recovered and how calculations are distributed among competing candidates.

The benchmark scores electronic energy differences. Zero-point, thermal/entropic, standard-state, and solvation contributions belong to a further free-energy treatment. This scope matters for NI5, specified in weakly acidic aqueous solution, and NU5, specified in dichloromethane at high concentration. Their scores assess the prescribed electronic-energy references. For association and dissociation, entropy and concentration-dependent standard-state terms enter $\Delta G^\ddagger$; solvent participation can also change the chemical model required to describe a pathway. Inferring solution-phase rate control or stereoselectivity would require a consistent free-energy treatment of the relevant species and pathways. The present electronic-energy comparisons leave the magnitude and direction of the solvent effect for a separate solvent-sensitivity study.

The atom-mapping coefficients and ambiguity window, and the decay, depth, and congestion settings in the branch-crowding descriptor, were fixed across the workflow comparisons. The study evaluates these fixed settings together; their sensitivity and individual contributions remain to be measured. The imaginary-mode alignment ratio of 0.5 is advisory; the stationary-point and IRC/endpoint evidence determine scientific validation. All P2 configurations used DeepSeek-v4.1-flash (1017), with the API alias \texttt{deepseek-v4-flash} and the recorded request settings. Comparisons with additional backbones would test how the workflow gains carry across models.

Extending Cloudra beyond small-molecule mechanism calculations will require representations, computational methods, and physical criteria appropriate to the new systems. Its existing computational engines and shared execution infrastructure can support those calculations. Cloudra is publicly available as open-source software under the GNU Affero General Public License (AGPL) at \url{https://github.com/yyywaa/Cloudra-mechanism-agent}, supporting inspection, reuse, and community development of the system. Reproducible tables and scores are in the companion data package.